% ===================================================================
%  اللوحة رقم ١٥ — السوق. --- المأكولات.
%
%  Traduction, faite en 2026, en arabe levantin d'aujourd'hui, sur
%  l'ido de texto/io. Regles de la colonne : voir 10-lawha-01.tex.
%  Une seule numerotation. 28 blocs, 102 renvois — le plus dense du
%  livret. L'ido coupe l'alinea 17 sur un « suite » entre les
%  feuillets 101 et 102 ; la traduction n'herite pas de la coupe et
%  ecrit UN bloc.
%
%  LA HALLE (1) EST « القيسريّة », ET LE MOT EST UN NOM D'EMPEREUR.
%  Le marche couvert de Royan recoit ici le nom que portent encore
%  les souks couverts d'Alep, de Damas et de Jerusalem : qaysariyya,
%  du grec καισάρεια, de Caesar. Le batiment romain a disparu, le mot
%  est reste accroche a la meme chose — un hall marchand sous toit —
%  et il est aujourd'hui sur les plaques de rue. C'EST LE TROISIEME
%  MOT DE CETTE COLONNE QUI EST SORTI DU LEVANT OU Y EST ENTRE PAR
%  L'EUROPE : المقهى au tableau 13, الترسانة et الأميرال au tableau
%  14, القيسريّة ici — sauf que celui-ci est entre et n'est jamais
%  ressorti.
%
%  TROIS FRUITS DANS UNE SEULE PHRASE, ET CHACUN DES TROIS DESIGNE
%  AUTRE CHOSE AU CAIRE. L'alinea 9 aligne la poire (52), la prune
%  (53) et la peche (55), ce qui n'arrive nulle part ailleurs dans le
%  livret. Le levantin ecrit إجّاص، خوخ، درّاق. L'egyptien du meme
%  releve ecrirait كمثرى، برقوق، خوخ. Les deux colonnes se servent
%  donc du mot خوخ, et il vaut prune de Beyrouth a Amman, peche du
%  Caire a Assouan. LA COLONNE VOISINE N'EST PAS UNE LANGUE
%  ETRANGERE, ET C'EST BIEN POUR CELA QUE LE PIEGE EST ICI : un mot
%  faux ne se verrait pas, il se lirait.
%
%  LES PETITS PAINS (82) SONT « الصمّون », ET C'EST LE DEUXIEME MOT
%  AKKADIEN DE CETTE COLONNE. L'ido ecrit « semli », qui est le
%  Semmel allemand ; le Semmel vient du latin simila, le latin du
%  grec σεμίδαλις, et le grec de l'akkadien samīdu — la fleur de
%  farine. Le meme mot akkadien est reste sur place par l'arameen et
%  donne aujourd'hui سميد la semoule et صمّون le petit pain. Le
%  tableau 5 avait deja trouve تموز, Tammuz, dans le calendrier
%  levantin. DEUX FOIS QUE CE RELEVE DESCEND SOUS L'ARABE, ET LES
%  DEUX FOIS PAR LA NOURRITURE ET PAR LE TEMPS.
%
%  ENFIN L'ALINEA 12 EST UN UNIFORME FRANCAIS DONT PAS UN MOT N'EST
%  FRANCAIS. Le dolman vient du turc dolama, le zouave du nom kabyle
%  Zwawa, la chechia de l'arabe شاشية — et le tarbouche (76) que
%  cette colonne ecrit طربوش porte le nom de Fes. Restent le hussard
%  et le shako, hongrois tous les deux. Rochelle habille deux soldats
%  de la Republique et pas un seul de leurs vetements n'a de nom
%  latin.
%
%  ET UN NEUVIEME METIER EN -جي : المكوجيّة la repasseuse (86), apres
%  les huit du tableau 14. La loi tient toujours — on ne repasse pas,
%  on tient une مكواة, et le metier prend le suffixe de l'objet.
% ===================================================================

%%K t15-apar-1 apar
\VUsaut{4.0mm}
\VUtitre{15.0pt}{60}{اللوحة رقم ١٥}
\VUsaut{3.0mm}

%%K t15-tit-1 sub
\VUpk{11.4pt}{40}{السوق. --- المأكولات.}
\VUsaut{3.0mm}

%%K t15-01-1 p
1. --- أكتر التجّار اللي بيجيبولنا الأكل اللي منحتاجه بيتجمّعوا تحت
\VUgras{القيسريّة} \textsuperscript{(1)} تبع \VUgras{السوق المسقوف}، أو
بالشوارع اللي حواليه.

%%K t15-02-1 p
2. --- \VUgras{لحّام الخنزير} \textsuperscript{(2)}، اللي بيبيع
\VUgras{الجامبون} \textsuperscript{(3)} و\VUgras{مقانق الدم}
\textsuperscript{(4)} و\VUgras{السجق} \textsuperscript{(5)}
و\VUgras{مقانق الكرشة} \textsuperscript{(6)} و\VUgras{شحم الخنزير}
\textsuperscript{(7)}، واقف جنب \VUgras{بيّاعة الزبدة والجبنة}
\textsuperscript{(8)}، اللي بسطتها معبّاية \VUgras{جبنة هولنديّة}
\textsuperscript{(9)} بقشرة حمرا، و\VUgras{جبنة كامامبير}
\textsuperscript{(10)}، و\VUgras{أقراص جبنة غرويير} \textsuperscript{(11)}
كبيرة، و\VUgras{قوالب زبدة} \textsuperscript{(12)} طازة.

%%K t15-03-1 p
3. --- وقدّامها، \VUgras{بيّاعة الخضرة} \textsuperscript{(13)} عم تعرض
على زبوناتها \VUgras{بندورة} \textsuperscript{(14)} حلوة،
و\VUgras{قرنبيط} \textsuperscript{(15)}، و\VUgras{خسّ}
\textsuperscript{(16)}، و\VUgras{ملفوف} \textsuperscript{(17)}،
و\VUgras{قرع} \textsuperscript{(19)}، و\VUgras{شمّام}
\textsuperscript{(18)}، وحتّى \VUgras{فطر} \textsuperscript{(20)}. بس
\VUgras{صبي البيرة}، اللي وقّف \VUgras{عرباية التوصيل}
\textsuperscript{(21)} تبعه وهي محمّلة \VUgras{براميل بيرة}
\textsuperscript{(22)} قدّام بسطتها بالزبط، عم يمنع زبوناتها يقرّبوا.
وفلّاحة عم تجبلها سلّة \VUgras{أرضي شوكي} \textsuperscript{(23)}.

%%K t15-tit-2 sub
\VUsaut{4.0mm}
\VUpk{10.2pt}{40}{بيع وشرا.}
\VUsaut{3.0mm}

%%K t15-04-1 p
4. --- وبالسوق الحركة كبيرة، لأنّه إجت ساعة \VUgras{المزاد}
\textsuperscript{(24)}. و\VUgras{ربّات البيوت} \textsuperscript{(26)}،
اللي جايين يتسوّقوا، بيوقفوا وهنّي ماشيين قدّام دكّان \VUgras{بيّاعة
الكرشة} \textsuperscript{(25)}، ليشتروا \VUgras{كرشة} أو \VUgras{راس
عجل.}

%%K t15-05-1 p
5. --- و\VUgras{طبّاخ} \textsuperscript{(27)} تخين عم يفاصل على
\VUgras{تونة} كتير حلوة عند \VUgras{بيّاعة السمك} \textsuperscript{(28)}
اللي بتغلّي أسعارها متل ما بدّها. إجاها إرساليّة كبيرة
\VUgras{إنكليس، وسمك موسى، وتروتة} و\VUgras{شبّوط.} و\VUgras{البكلا}
و\VUgras{بلح البحر} تبعها كتير طازة.

%%K t15-tit-3 sub
\VUsaut{4.0mm}
\VUpk{10.2pt}{40}{بضاعة رخيصة وبضاعة غالية.}
\VUsaut{3.0mm}

%%K t15-06-1 p
6. --- و\VUgras{بيّاعة الطيور} \textsuperscript{(29)}، اللي حاطّة
بسطتها جنبها، كتير أمينة. لمّا بتبيع \VUgras{ديك حبش، أو وزّ، أو بطّ}
أو \VUgras{جاجة} عاديّة، ما بتطلب غير السعر العادل.

%%K t15-07-1 p
7. --- وعند زاوية الساحة، عالطابق الأوّل، من مدّة قصيرة، انحطّ
\VUgras{بيّاع قماش} \textsuperscript{(30)} اللي عنده الشاريات دايمًا
أكيدين إنّهن بيلاقوا تشكيلة حلوة كتير من \VUgras{شراشف الطاولة، والفوط،
والمراييل} و\VUgras{مناشف الصحون.} وعلى نفس الطابق، \VUgras{سكّاكيني}
\textsuperscript{(31)} فتح \VUgras{معمله.} بيسنّ \VUgras{السكاكين} تبع
بيّاعات القيسريّة وبيعمل للفلّاحين \VUgras{مناجل} من أقسى
\VUgras{فولاذ.}

%%K t15-tit-4 sub
\VUsaut{4.0mm}
\VUpk{10.2pt}{40}{البقالة.}
\VUsaut{3.0mm}

%%K t15-08-1 p
8. --- وتحت معمله، من مدّة قصيرة، فتح محلّ مأكولات كبير. صار من زمان
مش هو \VUgras{دكّان البقالة} \textsuperscript{(32)} البسيط وقليل الشأن
تبع زمان، اللي كان بيبيع بس الأغراض اللي منستعملها كلّ يوم،
\VUgras{الشاي} \textsuperscript{(33)}، و\VUgras{الشوكولا}
\textsuperscript{(34)}، و\VUgras{قالب السكّر} \textsuperscript{(37)}
و\VUgras{الصابون} \textsuperscript{(38)}. هون بيبيعوا كلّ شي،
\VUgras{بسكوت،} و\VUgras{معلّبات} \textsuperscript{(35)}،
و\VUgras{مشروبات روحيّة} \textsuperscript{(36)}، و\VUgras{رَم، وكونياك،
وكيرش، وعرق يانسون} و\VUgras{كوراساو.} وبتلاقي هونيك صيد :
\VUgras{أرنب برّي} \textsuperscript{(39)}، و\VUgras{سمّن}
\textsuperscript{(40)}، و\VUgras{حجل} \textsuperscript{(41)}،
و\VUgras{فرّي} \textsuperscript{(42)}، و\VUgras{بطّ برّي}
\textsuperscript{(43)} أو \VUgras{تدرج} \textsuperscript{(44)}، متل ما
بتلاقي فواكه غريبة متل \VUgras{الموز} \textsuperscript{(46)}
و\VUgras{الأناناس.}

%%K t15-09-1 p
9. --- و\VUgras{صبي} \textsuperscript{(45)} البقّال، وهو كتير لطيف،
جاهز يقدّم اللي بدّك ياه : \VUgras{مربّى} \textsuperscript{(47)}
\VUgras{كشمش} أو \VUgras{سفرجل}، و\VUgras{شحم} \textsuperscript{(48)}،
و\VUgras{خيار مخلّل} \textsuperscript{(49)} أو \VUgras{سردين}
\textsuperscript{(50)} مملّح.

%%K t15-09-2 p
وهاد كتير مريح لـ\VUgras{الزبونات} \textsuperscript{(51)}، اللي
بيلاقوا، جنب أعدى الأغراض، أندر \VUgras{البواكير} وأطيب
\VUgras{الفواكه} : \VUgras{إجّاص} \textsuperscript{(52)} مسكّر،
و\VUgras{خوخ} \textsuperscript{(53)}، و\VUgras{عنب} \textsuperscript{(54)}،
و\VUgras{درّاق} \textsuperscript{(55)}، و\VUgras{لوز} \textsuperscript{(56)}
و\VUgras{جوز} \textsuperscript{(57)}.

%%K t15-tit-5 sub
\VUsaut{4.0mm}
\VUpk{10.2pt}{40}{الزباين.}
\VUsaut{3.0mm}

%%K t15-10-1 p
10. --- \VUgras{الرزّ} \textsuperscript{(58)} و\VUgras{العدس}
\textsuperscript{(59)} جنب صندوق \VUgras{الرنكة المدخّنة}
\textsuperscript{(60)} ؛ و\VUgras{البطاطا} \textsuperscript{(61)}
و\VUgras{الفاصوليا} \textsuperscript{(63)} جنب \VUgras{التفّاح}
\textsuperscript{(62)}، و\VUgras{التين} \textsuperscript{(64)}،
و\VUgras{القراصيا} \textsuperscript{(65)} و\VUgras{البندق}
\textsuperscript{(66)}. وبهالمحلّ بتلاقي \VUgras{بيض}
\textsuperscript{(67)} طازة و\VUgras{زيتون} \textsuperscript{(68)} ؛
لهيك \VUgras{السلال} \textsuperscript{(70)} و\VUgras{شبكات المونة}
\textsuperscript{(73)} بتتعبّى بلمح البصر.

%%K t15-11-1 p
11. --- و\VUgras{البنّ} \textsuperscript{(72)} تبع هالمحلّ الممتاز صار
إله سمعة حقيقيّة بين \VUgras{الشغّالات} \textsuperscript{(69)}
والطبّاخات، لأنّ \VUgras{البقّال} \textsuperscript{(71)} العتيق بيحمّصه
بإيده وبكلّ عناية.

%%K t15-tit-6 sub
\VUsaut{4.0mm}
\VUpk{10.2pt}{40}{المناظر.}
\VUsaut{3.0mm}

%%K t15-12-1 p
12. --- مش كلّ الناس بتيجي عالسوق لتتموّن. بالرايح والجاي المنظّر
تبع الزاروب اللي بيوصّل عالقيسريّة، بتشوف أغرب الناس. جنب
\VUgras{صبي المسلخ} \textsuperscript{(75)} الطيّب، اللي عم يدفع قدّامه
\VUgras{عرباية إيد} \textsuperscript{(74)}، هيّي عسكريّين بإجازة
مفتخرين كتير يفرجوا بدلتهن اللمّاعة : \VUgras{الهوسار}
\textsuperscript{(76)} بـ\VUgras{دولمان} أزرق و\VUgras{شاكو} بريش،
و\VUgras{عريف الزواف،} بشروال منفوخ ولابس \VUgras{طربوش.}

%%K t15-13-1 p
13. --- و\VUgras{الخبّاز} \textsuperscript{(80)}، اللي واقف على عتبة
\VUgras{الفرن} \textsuperscript{(78)}، لسّا عاجن عجينة \VUgras{خبزه}
\textsuperscript{(79)} وطالّع من التنّور \VUgras{الكرواسان}
\textsuperscript{(81)}، و\VUgras{الصمّون} \textsuperscript{(82)}
و\VUgras{أقراص الخبز} \textsuperscript{(83)} اللي بالفاترينة.

%%K t15-14-1 p
14. --- وفوق الفرن ساكنة \VUgras{خيّاطة بياضات} \textsuperscript{(84)}
شاطرة، عم تشغّل كذا \VUgras{مطرّزة} \textsuperscript{(85)}. وجنبها
ساكنة \VUgras{غسّالة ومكوجيّة} \textsuperscript{(86)}، اللي بتنشّي
\VUgras{البياضات} \textsuperscript{(88)} بعد ما تغسلها وتنشّفها،
وبتكويها بـ\VUgras{مكواة} \textsuperscript{(87)}.

%%K t15-tit-7 sub
\VUsaut{4.0mm}
\VUpk{10.2pt}{40}{لحمة وسمك وخضرة.}
\VUsaut{3.0mm}

%%K t15-15-1 p
15. --- وبـ\VUgras{الملحمة} \textsuperscript{(89)}، اللي عالطابق
الأرضي، بيبيعوا كلّ أنواع \VUgras{اللحمة، لحم غنم}
\textsuperscript{(90)}، و\VUgras{لحم بقر} \textsuperscript{(94)} أو
\VUgras{لحم عجل} \textsuperscript{(92)}، على شكل \VUgras{أفخاذ،
وبفتيك، ومشاوي} و\VUgras{ريش.} و\VUgras{صبي اللحّام}
\textsuperscript{(93)} بيفصّل كلّ هاللحمة وبيحطّ عالجنب \VUgras{عظام
النخاع} \textsuperscript{(96)}. وعلى باب الملحمة واقفة \VUgras{بيّاعة
المحار} \textsuperscript{(97)} اللي بتبيع \VUgras{محار}
\textsuperscript{(98)}. وبتفتحه إذا حدا بدّه ياه.

%%K t15-16-1 p
16. --- وكلّ هالتجّار بيدفعوا رخصة أو بدل مكان ؛ لهيك
\VUgras{الشرطي} \textsuperscript{(100)} بيلحق بلا رحمة على
\VUgras{بيّاعات الخضرة عالبسطات} \textsuperscript{(99)} المساكين اللي
بينافسوهن وهنّي شايلين بعربايتهن \VUgras{جزر، وفجل، ولفت، وكرّاث،} أو
\VUgras{ربطات هليون، وبازيلّا خضرا، وسبانخ، وحمّاض، ورشّاد}
و\VUgras{بقدونس.}

%%K t15-tit-8 sub
\VUsaut{4.0mm}
\VUpk{10.2pt}{40}{دير بالك من الشرطي !}
\VUsaut{3.0mm}

%%K t15-17-1 p
17. --- والولد اللي عم يبيع \VUgras{توم} \textsuperscript{(101)}
و\VUgras{بصل،} بيعرف منيح إنّه هو كمان ما بيحقّ له يبيع بضاعته جنب
السوق. بيهرب راكض، وناقصه يكبّ سلّة \VUgras{السرطانات، وجراد النهر}
و\VUgras{القريدس} تبع \VUgras{البيّاع} \textsuperscript{(102)} تبع
\VUgras{جراد البحر} و\VUgras{الكركند.}

%%K t15-18-1 p
18. --- والكلّ عم يتحرّك وعم يعمل ضجّة رهيبة ؛ بس هاد ما بيمنع إنّك
تسمع منيح صوت \VUgras{القزّاز} \textsuperscript{(103)} اللي عم يمرق
بالشارع، ولا صرخة \VUgras{الفحّام} \textsuperscript{(104)} اللي هلّق
ترك عربايته شوي ليوصّل \VUgras{كيس فحم} \textsuperscript{(105)}.
